\section{Related Work}
\label{sec:related}
% Sources: docs/related-work.md. Every claim here must be checked against the cited paper itself
% before submission; several entries in references.bib are still marked "unverified".

\paragraph{Learned pre-processing for standard codecs.}
Deep video precoding~\cite{bourtsoulatze2020deep} and deep perceptual
preprocessing~\cite{chadha2021deep} place a neural network in front of AVC, HEVC, VP9, AV1 or VVC
encoders and train it through a differentiable proxy of the codec; the decoder is unchanged.
Rate-perception optimised preprocessing~\cite{ma2023rate} is reported to run in a production
transcoding service. For user-generated content, a deep pre-processor tuned on pristine data was
found to make a substantial fraction of clips worse, which motivated adaptive control of the
pre-processing strength~\cite{zhao2025tridynamic}. Pre-editing before JPEG with a learned network
is the image analogue~\cite{talebi2021better}. All of these methods are applied by the party that
runs the encoder. Our setting differs in that the encoder, its configuration and the ladder that
chooses it are unknown to the optimiser. \todo{check each paper's exact setting and numbers.}

\paragraph{Training through non-differentiable codecs.}
Codec surrogates replace the real codec by a network during
training~\cite{tian2021self,chadha2021deep}. Alternatively, the real codec can be run in the
forward pass while a surrogate supplies the gradient~\cite{pesnel2026scaled}, an instance of the
straight-through and backward-pass-approximation
idea~\cite{bengio2013estimating,athalye2018obfuscated}. Accurate differentiable approximations exist
for simple codecs such as JPEG~\cite{reich2024differentiable}. Our emulator is a surrogate of the
entire platform pipeline rather than of a single codec.

\paragraph{Optimising against learned models.}
Optimising against a learned model drives the optimiser towards inputs on which the model
overestimates the objective. Offline model-based optimisation counters this with conservative
objective models~\cite{trabucco2021conservative}; model-based reinforcement learning penalises
ensemble disagreement~\cite{yu2020mopo}; trust-region Bayesian optimisation restricts the search
locally~\cite{eriksson2019scalable}. We observe the same failure when training a pre-compensation
through a platform emulator.

\paragraph{Metric hacking.}
VMAF~\cite{li2016toward} can be raised substantially by sharpening and by colour and contrast
changes without a subjective improvement~\cite{zvezdakova2019hacking}; the NEG variant reduces but
does not remove this~\cite{siniukov2021hacking,fingscheidt2026optimizing}. This is an instance of
Goodhart's law~\cite{manheim2018categorizing}. Our evaluation therefore separates the optimised
objective from judge metrics, following the recommendation to verify with metrics that were not
optimised.

\paragraph{Closed-loop adaptation to real platforms.}
Robust steganography uses transport channel matching: an image is uploaded to a social network
repeatedly until recompression no longer changes it~\cite{zhao2019improving}. This is the closest
methodological precedent for a closed loop against a real platform, with a different goal. Work on
the quality of user-generated content notes that the platform's reference is itself already
distorted~\cite{yu2021predicting,wang2019youtube}.

\paragraph{Black-box optimisation.}
We use Gaussian-process Bayesian optimisation~\cite{snoek2012practical,shahriari2016taking}, the
tree-structured Parzen estimator~\cite{bergstra2011algorithms}, CMA-ES~\cite{hansen2001completely}
and random search as implemented in Optuna~\cite{akiba2019optuna}.
